\documentclass[11pt]{article}
\usepackage[margin=1in]{geometry}
\usepackage{amsmath,amssymb,amsthm,mathtools}
\usepackage{hyperref}
\usepackage{enumitem}
\usepackage{xcolor}
\hypersetup{colorlinks=true, linkcolor=blue, urlcolor=blue, citecolor=blue}

\newtheorem{theorem}{Theorem}
\newtheorem{proposition}{Proposition}
\newtheorem{lemma}{Lemma}
\newtheorem{assumption}{Assumption}
\newtheorem{remark}{Remark}
\newtheorem{example}{Example}

\newcommand{\Q}{\mathcal Q}
\newcommand{\A}{\mathcal A}
\newcommand{\D}{\mathcal D}
\newcommand{\DeltaK}{\Delta_K}
\newcommand{\E}{\mathbb E}
\newcommand{\Pp}{\mathbb P}
\newcommand{\one}{\mathbf 1}
\newcommand{\what}{\widehat w}
\newcommand{\wstar}{w^\star}
\newcommand{\eps}{\varepsilon}
\newcommand{\norm}[1]{\left\lVert #1\right\rVert}

\title{A Two-Sample-Size Theory for Best-of-Infinity Ensembles:\\ Buffered ERM and Optional Fast Rates}
\author{Draft notes}
\date{\today}

\begin{document}
\maketitle

\section{Purpose of this note}

The goal is to isolate the statistical structure in Best-of-Infinity style LLM ensembling.  There are two independent sample sizes:
\[
\begin{aligned}
Q &= \text{number of labeled benchmark questions},\\
N &= \text{number of generations per model per question}.
\end{aligned}
\]
The first controls how well we learn ensemble weights across problems.  The second controls how well we estimate each model's answer distribution on each problem.  This note gives a clean theorem for a margin-buffered empirical risk minimizer and then states an optional fast-rate theorem under an additional geometric identifiability condition.

The main message is:
\[
\boxed{
R(\what)-R(\wstar)
\lesssim
\sqrt{\frac{K\log(AQ)+\log(1/\delta)}{Q}}
+
\left(\frac{\log(KAQ/\delta)}{N}\right)^{\beta/2}.}
\]
The $Q$-term is ordinary statistical generalization over labeled questions.  The $N$-term is the finite-generation floor.  The latter can be faster than $N^{-1/2}$ because the objective is a thresholded sign/correctness objective, not estimation of a smooth scalar parameter.

\section{Setup and notation}

Let $q\in\Q$ denote a problem, drawn from a distribution $\D$.  Each problem has a gold answer $a^\star(q)$ and a finite candidate answer set $\A(q)$ with
\[
|\A(q)|\le A.
\]
There are $K$ language models.  Model $i$ induces a distribution over final extracted answers,
\[
p_i(a\mid q),\qquad a\in\A(q),\quad i=1,\dots,K.
\]
For a fixed problem $q$, a fixed model $i$, and a generation budget $N$, we observe i.i.d. final answers
\[
Y_{i,1}(q),\dots,Y_{i,N}(q)\sim p_i(\cdot\mid q),
\]
and form the empirical answer distribution
\[
\widehat p_i(a\mid q)
:=
\frac1N\sum_{n=1}^N \one\{Y_{i,n}(q)=a\}.
\]
It is convenient to write the $K$-vector of model probabilities as
\[
p_\cdot(a\mid q)
:=
\big(p_1(a\mid q),\dots,p_K(a\mid q)\big)\in[0,1]^K,
\]
and similarly for $\widehat p_\cdot(a\mid q)$.

An ensemble is a weight vector
\[
w=(w_1,\dots,w_K)\in\Delta_K
:=
\left\{w\in\mathbb R_+^K:\sum_{i=1}^K w_i=1\right\}.
\]
Its true score for answer $a$ on problem $q$ is
\[
s_q(a;w):=\langle w,p_\cdot(a\mid q)\rangle,
\]
and its empirical score is
\[
\widehat s_q(a;w):=\langle w,\widehat p_\cdot(a\mid q)\rangle.
\]

For every wrong answer $a\ne a^\star(q)$ define the answer-gap vector
\[
\Delta_{q,a}
:=
p_\cdot(a^\star(q)\mid q)-p_\cdot(a\mid q)\in[-1,1]^K,
\]
and its empirical version
\[
\widehat\Delta_{q,a}
:=
\widehat p_\cdot(a^\star(q)\mid q)-\widehat p_\cdot(a\mid q).
\]
The true margin of $w$ on $q$ is
\[
M_q(w)
:=
\min_{a\in\A(q),\,a\ne a^\star(q)}
\langle w,\Delta_{q,a}\rangle.
\]
The empirical margin is
\[
\widehat M_q(w)
:=
\min_{a\in\A(q),\,a\ne a^\star(q)}
\langle w,\widehat\Delta_{q,a}\rangle.
\]
Thus $M_q(w)>0$ means that the true weighted ensemble assigns higher probability to the gold answer than to every wrong answer.  We count ties as errors.

The clean loss and clean risk are
\[
\ell_q(w):=\one\{M_q(w)\le 0\},
\qquad
R(w):=\Pp_{q\sim\D}\{M_q(w)\le 0\}.
\]
Let
\[
\wstar\in\arg\min_{w\in\Delta_K} R(w)
\]
be a clean-risk minimizer.

The training data are
\[
\mathcal S_{Q,N}
=
\left\{\left(q_j,a^\star(q_j),\{Y_{j,i,n}:i\in[K],n\in[N]\}\right):j=1,\dots,Q\right\},
\]
where $q_j\overset{i.i.d.}{\sim}\D$.  On the training set, $\widehat M_{q_j}(w)$ is observable because the benchmark supplies $a^\star(q_j)$ and the generations supply $\widehat p_i(\cdot\mid q_j)$.

\section{The margin-buffered estimator}

Raw ERM minimizes
\[
\frac1Q\sum_{j=1}^Q \one\{\widehat M_{q_j}(w)\le 0\}.
\]
This can exploit finite-$N$ noise near the decision boundary.  We instead use a buffer:
\[
\widehat R_\eps(w)
:=
\frac1Q\sum_{j=1}^Q \one\{\widehat M_{q_j}(w)\le \eps\},
\]
and define
\[
\what_\eps\in\arg\min_{w\in\Delta_K}\widehat R_\eps(w).
\]
The buffer is chosen at the scale of the uniform finite-generation error on the training sample:
\[
\eps_N(\delta)
:=
\sqrt{\frac{2\log(4KAQ/\delta)}{N}}.
\]
This choice is not sacred; any constant multiple of the same order works.

\section{A basic concentration event}

\begin{lemma}[Uniform training-sample margin accuracy]
\label{lem:margin-accuracy}
Fix $\delta\in(0,1)$.  With probability at least $1-\delta/2$ over the generation randomness conditional on $q_1,\dots,q_Q$,
\[
\sup_{j\le Q}\sup_{w\in\Delta_K}
\left|\widehat M_{q_j}(w)-M_{q_j}(w)\right|
\le
\eps_N(\delta).
\]
\end{lemma}

\begin{proof}
It suffices to control all answer frequencies.  By Hoeffding's inequality and a union bound,
\[
\Pp\left(
\max_{j\le Q}\max_{i\le K}\max_{a\in\A(q_j)}
\left|\widehat p_i(a\mid q_j)-p_i(a\mid q_j)\right|
>
\frac{\eps_N}{2}
\right)
\le
2KAQ\exp\left(-\frac{N\eps_N^2}{2}\right)
\le \frac\delta2.
\]
On this event, for every wrong answer $a$,
\[
\left|\langle w,\widehat\Delta_{q_j,a}-\Delta_{q_j,a}\rangle\right|
\le
\eps_N
\]
because $w\in\Delta_K$.  Taking minima over $a$ preserves the same Lipschitz bound.
\end{proof}

On the event in Lemma~\ref{lem:margin-accuracy}, for every training problem and every $w$,
\[
\one\{M_{q_j}(w)\le 0\}
\le
\one\{\widehat M_{q_j}(w)\le \eps_N\}
\le
\one\{M_{q_j}(w)\le 2\eps_N\}.
\]
This is the key sandwich.  The empirical noisy-buffered loss lies between the clean zero-buffer loss and the clean $2\eps_N$-buffer loss.

\section{Capacity of the clean margin classes}

For $t\ge0$ define the thresholded clean-margin class
\[
\mathcal F_t
:=
\left\{q\mapsto \one\{M_q(w)\le t\}:w\in\Delta_K\right\}.
\]
For fixed problems $q_1,
\dots,q_Q$, each pair $(q_j,a)$ with $a\ne a^\star(q_j)$ contributes the affine hyperplane
\[
\langle w,\Delta_{q_j,a}\rangle=t
\]
in the $(K-1)$-dimensional simplex.  There are at most $Q(A-1)$ such hyperplanes.  Therefore the number of possible sign patterns over $q_1,\dots,q_Q$ is at most the number of cells in an arrangement of $Q(A-1)$ hyperplanes in dimension $K-1$:
\[
\Pi_{\mathcal F_t}(Q)
\le
\sum_{r=0}^{K-1}\binom{Q(A-1)}{r}
\le
\left(\frac{eQ(A-1)}{K-1}\right)^{K-1}
\]
when $Q(A-1)\ge K-1$.  Thus the effective combinatorial dimension is of order $K\log(AK)$, or equivalently the finite-sample deviation carries a factor $(K-1)\log(AQ)$, not $KA$.

We will use the shorthand
\[
\mathfrak v_Q(\delta)
:=
C
\sqrt{
\frac{(K-1)\log(eAQ)+\log(1/\delta)}{Q}
},
\]
where $C>0$ is a universal constant.  Standard VC symmetrization and the above growth bound imply that, with probability at least $1-\delta/2$ over $q_1,\dots,q_Q$,
\[
\sup_{w\in\Delta_K,\,t\in\{0,2\eps_N\}}
\left|
\frac1Q\sum_{j=1}^Q\one\{M_{q_j}(w)\le t\}
-
\Pp_q\{M_q(w)\le t\}
\right|
\le
\mathfrak v_Q(\delta).
\]

\section{Main theorem: buffered ERM}

\begin{assumption}[Comparator margin condition]
\label{ass:comp-margin}
There are constants $C_m>0$, $\beta>0$, and $t_0>0$ such that
\[
\Pp_{q\sim\D}\{0<M_q(\wstar)\le t\}
\le
C_m t^\beta,
\qquad 0<t\le t_0.
\]
\end{assumption}

This is a one-sided condition: it only controls correctly solved problems at $\wstar$ whose positive margin is small.  That is enough for the buffered theorem because the proof compares the learned rule to $\wstar$ and only pays for the extra $2\eps_N$ band around correct decisions of $\wstar$.

\begin{theorem}[Buffered ERM, slow $Q$-rate and fast finite-generation floor]
\label{thm:buffered-main}
Assume Assumption~\ref{ass:comp-margin}.  Let
\[
\eps_N=\sqrt{\frac{2\log(8KAQ/\delta)}{N}}
\]
and suppose $2\eps_N\le t_0$.  Let
\[
\what_\eps\in\arg\min_{w\in\Delta_K}
\frac1Q\sum_{j=1}^Q\one\{\widehat M_{q_j}(w)\le\eps_N\}.
\]
Then with probability at least $1-\delta$ over both the $Q$ problem draws and the $N$ generations per model per problem,
\[
R(\what_\eps)-R(\wstar)
\le
2\mathfrak v_Q(\delta)
+
C_m(2\eps_N)^\beta.
\]
Equivalently, up to universal constants,
\[
\boxed{
R(\what_\eps)-R(\wstar)
\lesssim
\sqrt{\frac{(K-1)\log(eAQ)+\log(1/\delta)}{Q}}
+
C_m
\left(\frac{\log(KAQ/\delta)}{N}\right)^{\beta/2}.
}
\]
\end{theorem}

\begin{proof}
Let
\[
\widetilde R_t(w):=\frac1Q\sum_{j=1}^Q\one\{M_{q_j}(w)\le t\},
\qquad
R_t(w):=\Pp_q\{M_q(w)\le t\}.
\]
On the high-probability margin-accuracy event,
\[
\widetilde R_0(w)
\le
\widehat R_{\eps_N}(w)
\le
\widetilde R_{2\eps_N}(w)
\qquad
\text{for all }w\in\Delta_K.
\]
On the high-probability VC event,
\[
R_0(w)\le \widetilde R_0(w)+\mathfrak v_Q(\delta),
\qquad
\widetilde R_{2\eps_N}(w)\le R_{2\eps_N}(w)+\mathfrak v_Q(\delta)
\]
for all $w$.  Therefore
\begin{align*}
R(\what_\eps)
&=R_0(\what_\eps)\\
&\le \widetilde R_0(\what_\eps)+\mathfrak v_Q(\delta)\\
&\le \widehat R_{\eps_N}(\what_\eps)+\mathfrak v_Q(\delta)\\
&\le \widehat R_{\eps_N}(\wstar)+\mathfrak v_Q(\delta)\\
&\le \widetilde R_{2\eps_N}(\wstar)+\mathfrak v_Q(\delta)\\
&\le R_{2\eps_N}(\wstar)+2\mathfrak v_Q(\delta).
\end{align*}
Finally,
\[
R_{2\eps_N}(\wstar)-R_0(\wstar)
=
\Pp_q\{0<M_q(\wstar)\le 2\eps_N\}
\le
C_m(2\eps_N)^\beta
\]
by Assumption~\ref{ass:comp-margin}.  Subtract $R(\wstar)=R_0(\wstar)$.
\end{proof}

\begin{remark}[Why the buffer matters]
The theorem avoids any margin assumption at the random, data-dependent $\what_\eps$.  The buffer guarantees that if a training problem appears safely solved empirically, then it is also solved by the true margin, up to the $2\eps_N$ band.  Hence the only finite-$N$ mass we pay for is the $2\eps_N$ band around the fixed comparator $\wstar$.
\end{remark}

\section{Optional fast $Q$-rates from geometric identifiability}

The margin condition alone controls the finite-generation floor.  It does not, by itself, imply a fast $Q$-rate.  A fast $Q$-rate requires an additional condition saying that moving away from $\wstar$ actually increases risk.  This is an identifiability or growth condition.

Define the disagreement probability
\[
D(w,\wstar)
:=
\Pp_q\{\ell_q(w)\ne \ell_q(\wstar)\}.
\]
A common route to fast rates is a Bernstein condition
\[
D(w,\wstar)
\le
B\{R(w)-R(\wstar)\}^{\alpha},
\qquad 0<\alpha\le 1.
\]
Rather than assume this directly, we state primitive geometric conditions that imply it.

\begin{assumption}[Absolute comparator margin]
\label{ass:absolute-margin}
There exist $C_m>0$, $\beta>0$, and $t_0>0$ such that
\[
\Pp_q\{|M_q(\wstar)|\le t\}
\le
C_m t^\beta,
\qquad 0<t\le t_0.
\]
\end{assumption}

\begin{assumption}[Local risk growth and global separation]
\label{ass:growth}
There exist constants $c_g>0$, $\gamma>0$, $r_0>0$, and $c_{\rm sep}>0$ such that
\[
R(w)-R(\wstar)
\ge
c_g\norm{w-\wstar}_1^\gamma
\qquad
\text{whenever }\norm{w-\wstar}_1\le r_0,
\]
and
\[
R(w)-R(\wstar)
\ge c_{\rm sep}
\qquad
\text{whenever }\norm{w-\wstar}_1>r_0.
\]
\end{assumption}

\begin{proposition}[Geometry implies Bernstein]
\label{prop:geom-bernstein}
Under Assumptions~\ref{ass:absolute-margin} and~\ref{ass:growth}, there is a constant $B<\infty$ such that
\[
D(w,\wstar)
\le
B\{R(w)-R(\wstar)\}^{\alpha}
\qquad
\text{for all }w\in\Delta_K,
\]
where
\[
\alpha:=\min\left\{1,\frac{\beta}{\gamma}\right\}.
\]
\end{proposition}

\begin{proof}
First note that $M_q$ is one-Lipschitz in $\ell_1$:
\[
|M_q(w)-M_q(w')|
\le
\norm{w-w'}_1.
\]
Indeed, each $\Delta_{q,a}$ has $\ell_\infty$ norm at most one, and taking the minimum over $a$ preserves the Lipschitz constant.  If $\ell_q(w)\ne\ell_q(\wstar)$, then the signs of $M_q(w)$ and $M_q(\wstar)$ differ, so
\[
|M_q(\wstar)|
\le
|M_q(w)-M_q(\wstar)|
\le
\norm{w-\wstar}_1.
\]
Thus, by Assumption~\ref{ass:absolute-margin},
\[
D(w,\wstar)
\le
C_m\norm{w-\wstar}_1^\beta
\]
whenever $\norm{w-\wstar}_1\le t_0$.  Locally, Assumption~\ref{ass:growth} gives
\[
\norm{w-\wstar}_1^\beta
\le
c_g^{-\beta/\gamma}
\{R(w)-R(\wstar)\}^{\beta/\gamma}.
\]
If $\beta/\gamma>1$, weaken the exponent to one, since $R(w)-R(\wstar)\le 1$ after adjusting constants.  Away from the local ball, $D(w,\wstar)\le1\le c_{\rm sep}^{-\alpha}\{R(w)-R(\wstar)\}^{\alpha}$.  Enlarging constants proves the claim.
\end{proof}

We will use the following standard localized-VC lemma.  It is included as a black-box empirical-process step; the substantive modeling content for this paper is Proposition~\ref{prop:geom-bernstein}, which gives primitive geometric conditions under which the Bernstein premise of the lemma holds.

\begin{lemma}[Localized VC bound for approximate ERM]
\label{lem:localized-vc}
Let $\mathcal L=\{\ell_w:w\in\Delta_K\}$ be a class of $\{0,1\}$-valued losses with minimizer $\wstar$.  Suppose its growth function on $Q$ points is bounded by
\[
\Pi_{\mathcal L}(Q)\le \exp(V_Q),
\qquad
V_Q\asymp (K-1)\log(eAQ),
\]
and suppose the Bernstein condition holds:
\[
\Pp_q\{\ell_q(w)\ne\ell_q(\wstar)\}
\le
B\{R(w)-R(\wstar)\}^{\alpha},
\qquad 0<\alpha\le 1.
\]
If a data-dependent $\widetilde w$ satisfies
\[
\frac1Q\sum_{j=1}^Q\ell_{q_j}(\widetilde w)
\le
\frac1Q\sum_{j=1}^Q\ell_{q_j}(\wstar)+a,
\]
then with probability at least $1-\delta$,
\[
R(\widetilde w)-R(\wstar)
\le
C
\left[
 a
 +
\left(\frac{V_Q+\log(1/\delta)}{Q}\right)^{1/(2-\alpha)}
+
\frac{V_Q+\log(1/\delta)}{Q}
\right],
\]
where $C$ depends only on $B$ and $\alpha$.
\end{lemma}

\begin{proof}[Proof idea]
This is the usual peeling proof for bounded excess losses under a Bernstein condition.  For the excess class $g_w=\ell_w-\ell_{\wstar}$, the variance satisfies $\E g_w^2\le B(\E g_w)^\alpha$.  On each shell $r<\E g_w\le2r$, Bernstein's inequality and the VC growth bound imply a uniform deviation of order
\[
\sqrt{\frac{r^\alpha(V_Q+\log(1/\delta))}{Q}}
+
\frac{V_Q+\log(1/\delta)}{Q}.
\]
Solving the fixed point
\[
r\asymp
\sqrt{\frac{r^\alpha(V_Q+\log(1/\delta))}{Q}}
\]
gives $r\asymp((V_Q+\log(1/\delta))/Q)^{1/(2-\alpha)}$.  The approximate-ERM slack $a$ enters additively.  This is the standard Massart/Tsybakov localized-rate argument, specialized to the finite-growth class above.
\end{proof}

We now state the fast-rate version.

\begin{theorem}[Buffered ERM with optional fast $Q$-rate]
\label{thm:fast}
Assume Assumptions~\ref{ass:absolute-margin} and~\ref{ass:growth}.  Let $\alpha=\min\{1,\beta/\gamma\}$ and let $\what_\eps$ be the buffered ERM from Theorem~\ref{thm:buffered-main}.  Then, with probability at least $1-\delta$,
\[
R(\what_\eps)-R(\wstar)
\le
C
\left(
\frac{(K-1)\log(eAQ)+\log(1/\delta)}{Q}
\right)^{\frac{1}{2-\alpha}}
+
C'
\left(\frac{\log(KAQ/\delta)}{N}\right)^{\beta/2},
\]
where $C,C'$ depend on the margin and growth constants but not on $Q,N,K,A$.
\end{theorem}

\begin{proof}[Proof sketch]
On the margin-accuracy event,
\[
\frac1Q\sum_{j=1}^Q\ell_{q_j}(\what_\eps)
\le
\frac1Q\sum_{j=1}^Q\one\{\widehat M_{q_j}(\what_\eps)\le\eps_N\}
\le
\frac1Q\sum_{j=1}^Q\one\{\widehat M_{q_j}(\wstar)\le\eps_N\}
\le
\frac1Q\sum_{j=1}^Q\one\{M_{q_j}(\wstar)\le2\eps_N\}.
\]
Therefore $\what_\eps$ is an approximate ERM for the clean loss, with empirical excess bounded by
\[
\frac1Q\sum_{j=1}^Q\one\{0<M_{q_j}(\wstar)\le2\eps_N\}.
\]
By the margin condition, this empirical band mass has expectation at most $C_m(2\eps_N)^\beta$ and, by Bernstein's inequality for a fixed Bernoulli variable, is at most a constant multiple of $(2\eps_N)^\beta+\log(1/\delta)/Q$ with high probability.  Proposition~\ref{prop:geom-bernstein} gives the Bernstein condition with exponent $\alpha$.  Applying Lemma~\ref{lem:localized-vc} yields the displayed bound, absorbing the $\log(1/\delta)/Q$ term into the localized complexity term.
\end{proof}

\begin{remark}[Where the exponent comes from]
The exponent is not inherently $\beta/(\beta+1)$.  In this geometry it is
\[
\alpha=\min\{1,\beta/\gamma\},
\]
where $\beta$ measures the thickness of the decision boundary and $\gamma$ measures how risk grows as one moves away from $\wstar$.  The often-appearing value $\alpha=\beta/(\beta+1)$ corresponds to the special local-growth regime $\gamma=\beta+1$.
\end{remark}

\section{Why margin alone cannot give a fast $Q$-rate}

The following example shows why the fast $Q$-rate cannot follow from the geometric margin condition alone, even with a hard margin and no finite-generation noise.

\begin{proposition}[Hard margin but only root-$Q$ distinguishability]
\label{prop:counterexample}
There is a family of two-model, two-answer problems with $N=\infty$ and hard margin at the optimum such that no procedure can uniformly guarantee excess risk $o(Q^{-1/2})$ from $Q$ labeled questions.
\end{proposition}

\begin{proof}
Take $K=2$ and two answers, the gold answer $a^\star$ and a wrong answer $b$.  Write $w=(t,1-t)$.  There are two problem types.

For type $+$, model 1 is always correct and model 2 is always wrong:
\[
p_1(a^\star\mid q)=1,
\qquad
p_2(a^\star\mid q)=0.
\]
Then
\[
\Delta_q=(1,-1),
\qquad
M_q(w)=2t-1.
\]
For type $-$, model 1 is always wrong and model 2 is always correct:
\[
p_1(a^\star\mid q)=0,
\qquad
p_2(a^\star\mid q)=1,
\]
so
\[
\Delta_q=(-1,1),
\qquad
M_q(w)=1-2t.
\]
Thus weights with $t>1/2$ solve exactly the type $+$ problems, while weights with $t<1/2$ solve exactly the type $-$ problems.

Now consider two distributions over problem types:
\[
\D_+:
\Pp(q=+)=\frac12+\frac\eps2,
\qquad
\D_-:
\Pp(q=+)=\frac12-\frac\eps2.
\]
Under $\D_+$, the optimal side is $t>1/2$, with risk $1/2-\eps/2$.  Choosing the wrong side has risk $1/2+\eps/2$, so the excess risk is $\eps$.  Under $\D_-$, the optimal side is reversed.

At an optimum, say $w^\star=(1,0)$ under $\D_+$, the margins satisfy
\[
|M_q(w^\star)|=1
\]
for both problem types.  Hence the margin is hard: there is no near-boundary mass at all.

However, distinguishing $\D_+$ from $\D_-$ using $Q$ training questions is just distinguishing two coins with biases $1/2\pm\eps/2$.  If $\eps=c/\sqrt Q$ for sufficiently small constant $c$, the total variation distance between the $Q$-sample laws remains bounded away from one.  Therefore any procedure chooses the wrong side with constant probability under at least one of $\D_+$ or $\D_-$.  On that event its excess risk is $\eps\asymp Q^{-1/2}$.

Thus even a hard margin does not imply a faster-than-root-$Q$ learning rate.  The missing ingredient is global identifiability: the two separated regions of the simplex have nearly equal risks.
\end{proof}

\section{How to interpret the growth exponent $\gamma$}

The growth condition is not a technical decoration; it describes how the geometry of the benchmark changes when we move the ensemble weights.  Let $r=\norm{w-\wstar}_1$.  Because $M_q$ is Lipschitz in $w$, only problems with
\[
|M_q(\wstar)|\lesssim r
\]
can change their solved/unsolved status.  Under the margin condition, the mass of such problems is at most order $r^\beta$.  This is the disagreement scale:
\[
D(w,\wstar)\lesssim r^\beta.
\]
But the excess risk is not the disagreement mass; it is the net imbalance between newly lost and newly gained questions.

Along a one-dimensional path $w_r$ away from $\wstar$, write
\[
B_+(r):=\Pp\{\text{question is solved by }\wstar\text{ but failed by }w_r\},
\]
\[
B_-(r):=\Pp\{\text{question is failed by }\wstar\text{ but solved by }w_r\}.
\]
Then
\[
D(w_r,\wstar)=B_+(r)+B_-(r),
\qquad
R(w_r)-R(\wstar)=B_+(r)-B_-(r).
\]
The margin exponent controls the sum; the growth exponent controls the difference.

\paragraph{Case A: all crossings hurt.}
If moving away from $\wstar$ mostly turns correct questions into errors and does not rescue comparable failed questions, then
\[
B_+(r)\asymp r^\beta,
\qquad
B_-(r)=o(r^\beta).
\]
Hence
\[
R(w_r)-R(\wstar)\asymp r^\beta,
\]
so $\gamma=\beta$, $\alpha=1$, and the fast $Q$-term is essentially $Q^{-1}$ up to logarithms.  A simple model is a one-sided boundary: all questions near the boundary are correctly solved at $\wstar$ and become wrong when the weight is perturbed.

\paragraph{Case B: first-order cancellation.}
Suppose the total crossing mass is still order $r^\beta$, but losses and gains nearly cancel:
\[
B_+(r)=c r^\beta+c' r^{\beta+1}+o(r^{\beta+1}),
\]
\[
B_-(r)=c r^\beta-c' r^{\beta+1}+o(r^{\beta+1}).
\]
Then
\[
D(w_r,\wstar)\asymp r^\beta,
\qquad
R(w_r)-R(\wstar)\asymp r^{\beta+1}.
\]
Thus $\gamma=\beta+1$, $\alpha=\beta/(\beta+1)$, and the fast $Q$-term becomes
\[
Q^{-(\beta+1)/(\beta+2)}
\]
up to logarithms.  This is the regime in which the previously suggested exponent naturally appears.  It is not a consequence of the margin condition alone; it is a consequence of margin plus first-order gain/loss cancellation.

\paragraph{Case C: higher-order cancellation.}
If the gain/loss imbalance vanishes to higher order,
\[
B_+(r)-B_-(r)\asymp r^{\beta+s},
\qquad s>1,
\]
then $\gamma=\beta+s$ and
\[
\alpha=\frac{\beta}{\beta+s}.
\]
The fast rate becomes slower.  This is a flatter risk landscape: many questions change status near the boundary, but the net risk difference is small.

\paragraph{Case D: separated near-tie between two regions.}
The counterexample in Proposition~\ref{prop:counterexample} is the extreme nonlocal case.  Two separated regions of the simplex have hard margins internally, but their risks differ by only $\eps$.  The benchmark gives no local boundary signal that decides which region is better; one must estimate a global coin bias from $Q$ problems.  This yields the root-$Q$ obstruction.

\section{Recommended theorem hierarchy for the paper}

The cleanest presentation is:

\begin{enumerate}[label=(\roman*)]
\item State the double-randomness setup and the definitions of $p_i$, $\widehat p_i$, $\Delta_{q,a}$, $M_q$, and $\widehat M_q$.

\item Make the buffered ERM theorem, Theorem~\ref{thm:buffered-main}, the main result.  It requires only a comparator margin condition and gives the robust two-term rate
\[
\sqrt{\frac{K\log(AQ)}{Q}}
+
\left(\frac{\log(KAQ)}{N}\right)^{\beta/2}.
\]

\item Include Proposition~\ref{prop:counterexample} immediately after the main theorem to explain why the margin condition alone cannot produce a fast $Q$-rate.

\item Present the optional fast-rate theorem only after introducing geometric identifiability: absolute margin plus risk growth.  Emphasize that the exponent is
\[
\alpha=\min\{1,\beta/\gamma\},
\]
not automatically $\beta/(\beta+1)$.
\end{enumerate}

This makes the paper honest: the genuinely new LLM-specific contribution is the $Q,N$ decomposition and the finite-generation margin floor.  Fast $Q$-rates are available, but only when the risk landscape has additional identifiable geometry.

\end{document}
